% Canonical theory homes C10.01--C10.14. No solved-question/revision duplicates.
\begingroup
\chTwoReadableTopicStyle
\widowpenalty=10000\clubpenalty=10000
\section{ক্লাউড কম্পিউটিং}
\subsection{Cloud Computing: কেন প্রয়োজন, কীভাবে কাজ করে}
\label{ch2:c10-01}
\begin{definitionbox}{Cloud Computing}
নেটওয়ার্কের মাধ্যমে ভাগাভাগি করা ও প্রয়োজনমতো বিন্যস্ত করা যায় এমন কম্পিউটিং সম্পদ থেকে চাহিদা অনুযায়ী দ্রুত সার্ভার, স্টোরেজ, নেটওয়ার্ক ও অ্যাপ্লিকেশন নেওয়া এবং ছেড়ে দেওয়ার ব্যবস্থাকে ক্লাউড কম্পিউটিং বলে। এতে প্রতিবার সেবাদাতার কর্মীর সঙ্গে যোগাযোগ করে ব্যবস্থা নেওয়ার প্রয়োজন কম থাকে।
\end{definitionbox}
কলেজের ভর্তি ওয়েবসাইটে কয়েক দিনের জন্য অনেক ব্যবহারকারী আসে। সারা বছরের জন্য সর্বোচ্চ ক্ষমতার সার্ভার কিনলে অধিকাংশ সময় তার ব্যবহার কম হতে পারে। ক্লাউডে প্রয়োজন অনুযায়ী সম্পদ নেওয়া যায়। অনলাইন স্টোরেজের পাশাপাশি অ্যাপ্লিকেশন, ডেটাবেজ ও প্রক্রিয়াকরণও ক্লাউড সেবা হতে পারে।
\chTwoPair{\textbf{ব্যবস্থার উপাদান ও কাজ}\par ফোন বা ল্যাপটপ ক্লায়েন্ট হিসেবে নেটওয়ার্ক দিয়ে অনুরোধ পাঠায়। সার্ভার অ্যাকাউন্ট ও অনুমতি যাচাই করে অ্যাপ্লিকেশন চালায়। স্টোরেজ বা ডেটাবেজ থেকে তথ্য নেয় কিংবা সেখানে তথ্য রাখে এবং ফল ক্লায়েন্টে ফেরত পাঠায়। সেবাদাতার ডেটা সেন্টারে ভৌত সার্ভার, নেটওয়ার্ক, স্টোরেজ ও সেবা পরিচালনার সফটওয়্যার থাকে।}{\centering
\begin{tikzpicture}[node distance=8mm]
\node[chTwoNode,text width=2.8cm] (u) {Client\\Phone/Laptop};
\node[chTwoNode,text width=2.8cm,below=of u] (n) {Network};
\node[chTwoNode,text width=2.8cm,below=of n] (s) {Cloud Service\\Server + Storage};
\draw[chTwoLink,<->] (u)--(n);
\draw[chTwoLink,<->] (n)--(s);
\end{tikzpicture}\par Request যায়; Data/ফল ফিরে আসে।}
Public Cloud-এর সেবা সাধারণত ইন্টারনেট দিয়ে পাওয়া যায়। Private Cloud নিজস্ব নেটওয়ার্কেও চলতে পারে। দূরের কম্পিউটারে ফাইল রাখা বা ওয়েবসাইট পরিচালনা করলেই ক্লাউডের সব বৈশিষ্ট্য পূরণ হয় না। \textbf{Virtualization}-এ একটি ভৌত যন্ত্রের সম্পদ দিয়ে একাধিক ভার্চুয়াল যন্ত্র চলে। এটি ক্লাউড তৈরিতে সাহায্য করে। তবে ক্লাউডে সম্পদ বরাদ্দ, প্রবেশাধিকার, ব্যবস্থাপনা ও ব্যবহার মাপার ব্যবস্থাও প্রয়োজন।

\subsection{পাঁচটি মৌলিক বৈশিষ্ট্য}
\label{ch2:c10-02}
ক্লাউডের ব্যবহার বোঝার তিনটি সেবানীতি হলো: \textbf{চাহিদামতো সেবা নেওয়া}, \textbf{চাহিদা অনুযায়ী সম্পদ বাড়ানো বা কমানো} এবং \textbf{ব্যবহারের পরিমাণ অনুযায়ী মূল্য দেওয়া}। এগুলো সুবিধা বোঝানোর সংক্ষিপ্ত ছক। পূর্ণ মান অনুযায়ী বৈশিষ্ট্যের তালিকা নিচে দেওয়া হয়েছে।

NIST (National Institute of Standards and Technology)-এর ২০১১ সালের SP 800-145 সংজ্ঞায় পাঁচটি মৌলিক বৈশিষ্ট্য, তিনটি সেবা মডেল ও চারটি স্থাপনা মডেল রয়েছে।
\begin{description}
\item[On-demand Self-service:] ব্যবহারকারী চাহিদা অনুযায়ী নিজেই স্টোরেজ বা প্রক্রিয়াকরণের সম্পদ চেয়ে নিতে পারেন। প্রতিবার সেবাদাতার কর্মীর সঙ্গে যোগাযোগ প্রয়োজন হয় না।
\item[Broad Network Access:] নেটওয়ার্কের প্রচলিত ব্যবস্থায় বিভিন্ন ক্লায়েন্ট ডিভাইস থেকে সেবা ব্যবহার করা যায়।
\item[Resource Pooling:] সেবাদাতার ভৌত ও ভার্চুয়াল সম্পদের একটি ভান্ডার থেকে ব্যবহারকারীর চাহিদা অনুযায়ী সম্পদ বরাদ্দ ও পুনর্বণ্টন করা হয়। গ্রাহক সাধারণত নির্দিষ্ট ভৌত সার্ভার নির্বাচন করেন না।
\item[Rapid Elasticity:] চাহিদা বাড়লে দ্রুত সম্পদ বাড়ানো এবং কমলে তা ছেড়ে দেওয়া যায়। ভর্তি চলাকালে সক্ষমতা বাড়িয়ে পরে কমানো এর উদাহরণ।
\item[Measured Service:] স্টোরেজ, প্রক্রিয়াকরণ, ডেটা স্থানান্তর বা অ্যাকাউন্টের ব্যবহার মাপা, নিয়ন্ত্রণ ও প্রতিবেদন করা যায়। ব্যবহারের সঙ্গে বিল সম্পর্কিত হতে পারে। সব সেবার মূল্য নির্ধারণের পদ্ধতি একই নয়।
\end{description}

\subsection{Service Model: কী পাওয়া যায়, কে পরিচালনা করে}
\label{ch2:c10-03}
\begin{tabularx}{\linewidth}{|p{.18\linewidth}|X|X|}\hline
\textbf{Model} & \textbf{সেবা ও ব্যবহারকারীর দায়িত্ব} & \textbf{ব্যবহারের উদাহরণ}\\\hline
IaaS: Infrastructure as a Service & ভার্চুয়াল যন্ত্র, স্টোরেজ ও নেটওয়ার্ক পাওয়া যায়। সেবাদাতা ভৌত অবকাঠামো চালায়। ব্যবহারকারী সাধারণত OS, অ্যাপ্লিকেশন ও ডেটা পরিচালনা করেন। & ভর্তি, পরীক্ষা বা ব্যাকআপের সার্ভার; যেমন AWS EC2, Azure Virtual Machines, Google Compute Engine। OS ও অ্যাপ্লিকেশন হালনাগাদ এবং কনফিগারেশন নিজে করতে হয়।\\\hline
PaaS: Platform as a Service\label{ch2:c10-04} & সেবাদাতার প্রস্তুত রানটাইম ও উন্নয়ন প্ল্যাটফর্মে অ্যাপ্লিকেশন তৈরি ও চালানো যায়। ব্যবহারকারী অ্যাপ্লিকেশন ও ডেটা পরিচালনা করেন। OS ও রানটাইম ব্যবস্থাপনার বড় অংশ সেবাদাতা করে। & ওয়েবসাইটের কোড স্থাপন; যেমন Google App Engine, Azure App Service, Heroku। সেবাদাতা OS ও প্ল্যাটফর্ম চালায়। অ্যাপ্লিকেশনকে সমর্থিত ভাষা, টুল ও সীমার সঙ্গে মানিয়ে নিতে হয়।\\\hline
SaaS: Software as a Service\label{ch2:c10-05} & প্রস্তুত অ্যাপ্লিকেশন ব্যবহার করা যায়। সেবাদাতা কেন্দ্রীয় সার্ভার ও অ্যাপ্লিকেশন হালনাগাদ করে। ব্যবহারকারীকে নিজের সার্ভারে প্রতিবার নতুন সংস্করণ বসাতে হয় না। নিজের ডেটা, অ্যাকাউন্ট ও অনুমতি পরিচালনা করতে হয়। & Gmail, Google Docs/Sheets, Microsoft 365, Dropbox, Canva বা Zoom ব্যবহার। নিজস্ব সার্ভার বা অ্যাপ্লিকেশন তৈরি করতে হয় না।\\\hline
\end{tabularx}
IaaS থেকে SaaS-এর দিকে ব্যবহারকারীর পরিচালনাগত দায়িত্ব কমে। তবে সবক্ষেত্রে অ্যাকাউন্ট ও ডেটার নিয়ন্ত্রণ এবং নিরাপদ ব্যবহার প্রয়োজন। SaaS ব্যবহার করলেও নিজের ডিভাইসের OS ও ব্রাউজার হালনাগাদ রাখতে হয়। একই সেবাদাতা বিভিন্ন মডেলের সেবা দিতে পারে। নির্দিষ্ট সেবার কাজ দেখে মডেল শনাক্ত করতে হয়।
\subsection{Deployment Model: কার জন্য Infrastructure}
\label{ch2:c10-06}
\begin{longtable}{|p{\dimexpr .18\linewidth-7pt\relax}|p{\dimexpr .40\linewidth-7pt\relax}|p{\dimexpr .42\linewidth-7pt\relax}|}\hline
\textbf{Model} & \textbf{ব্যবহার ও নিয়ন্ত্রণ} & \textbf{উদাহরণ, সুবিধা ও সীমা}\\\hline
\endfirsthead\hline\textbf{Model} & \textbf{ব্যবহার ও নিয়ন্ত্রণ} & \textbf{উদাহরণ, সুবিধা ও সীমা}\\\hline\endhead
Public & সাধারণ গ্রাহকের জন্য সেবাদাতা সম্পদ পরিচালনা করে। সেবাদাতা ব্যবসায়িক, সরকারি বা শিক্ষাপ্রতিষ্ঠান হতে পারে। & বড় প্রাথমিক হার্ডওয়্যার ব্যয় ছাড়াই সেবা পাওয়া যায়। Public মানে সব ফাইল উন্মুক্ত বা সব সেবা বিনামূল্যে নয়।\\\hline
Private\label{ch2:c10-07} & এক প্রতিষ্ঠানের একক ব্যবহারের ব্যবস্থা। নিজে বা অন্য সেবাদাতা পরিচালনা করতে পারে। প্রতিষ্ঠানের ভেতরে বা বাইরে থাকতে পারে। & ব্যাংকের নিজস্ব ব্যবস্থাপনা বা কলেজের সংবেদনশীল ফলাফল। নিয়ন্ত্রণ ও নিজস্ব প্রয়োজন অনুযায়ী বিন্যাসের সুযোগ বেশি। দক্ষ দল ও অর্থ প্রয়োজন। নিরাপত্তার ঝুঁকি শূন্য হয় না।\\\hline
Community\label{ch2:c10-08} & অভিন্ন উদ্দেশ্য, নীতি বা নিরাপত্তার প্রয়োজন রয়েছে এমন কয়েকটি প্রতিষ্ঠানের নির্দিষ্ট গোষ্ঠীর জন্য। & হাসপাতাল, গবেষণা বা সরকারি প্রতিষ্ঠানের যৌথ ব্যবস্থা। খরচ ভাগ করা যায়; দায়িত্ব ও নীতিতে সমঝোতা প্রয়োজন। সাধারণ Public অ্যাকাউন্ট থাকা এর প্রমাণ নয়।\\\hline
Hybrid\label{ch2:c10-09} & পৃথক Public/Private/Community Cloud নিজস্ব পরিচয় বজায় রেখে ডেটা বা অ্যাপ্লিকেশনের জন্য প্রযুক্তিগতভাবে সংযুক্ত থাকে। & ব্যক্তিগত ফলাফল ব্যবস্থা ও সাধারণের ভর্তি সেবার সমন্বয়। নিয়ন্ত্রণ ও চাহিদা সামলানো যায়; সমন্বয় ও প্রবেশাধিকার পরিচালনা জটিল। স্থানীয় কম্পিউটার ও পৃথক Drive পাশাপাশি ব্যবহার যথেষ্ট নয়।\\\hline
\end{longtable}
\subsection{Service ও Deployment: দুই স্বতন্ত্র ভিত্তি}
\label{ch2:c10-10}
\chTwoPair{\textbf{Service Model}\par কী সম্পদ পাওয়া যায় এবং কোন স্তর নিজে পরিচালনা করতে হয়, তা বোঝায়। IaaS/PaaS/SaaS-এর মূল ব্যাখ্যা অংশ \ref{ch2:c10-03}-এ আছে।}{\textbf{Deployment Model}\par অবকাঠামো কার ব্যবহারের জন্য এবং কীভাবে নিয়ন্ত্রিত বা সমন্বিত, তা বোঝায়। Public/Private/Community/Hybrid-এর ব্যাখ্যা অংশ \ref{ch2:c10-06}-এ আছে।}
প্রস্তুত অনলাইন ডকুমেন্ট সেবা একই সঙ্গে SaaS ও Public Deployment হতে পারে। পরিচয় দুটি পরস্পরের বিকল্প নয়। ফাইল রাখার সুবিধা অনুযায়ী Google Drive ক্লাউড স্টোরেজ এবং প্রস্তুত অ্যাপ্লিকেশন হিসেবে SaaS। Azure/AWS সেবাদাতা বা সেবাসমষ্টির নাম। তাদের ভার্চুয়াল যন্ত্র ও পরিচালিত অ্যাপ্লিকেশন প্ল্যাটফর্মের মডেল আলাদা হতে পারে। শুধু সেবাদাতার নাম দিয়ে নির্দিষ্ট সেবা মডেল নির্ধারণ করা যায় না।
\chTwoPair{\textbf{NaaS (Network as a Service):} গ্রাহকচুক্তিতে নেটওয়ার্ক সংযোগ, ব্যান্ডউইথ, VPN ও সংশ্লিষ্ট ব্যবস্থাপনা সেবা পাওয়া যায়। যেমন, শাখার সংযোগের সব যন্ত্র নিজে পরিচালনা না করে সেবাদাতার কাছ থেকে সেবা নেওয়া। এখানে নেটওয়ার্কের সম্পদ ও সেবা দেওয়া মূল বিষয়; প্রস্তুত সফটওয়্যার ব্যবহার নয়।}{\textbf{CuaS: ক্লায়েন্ট দিয়ে ক্লাউডে প্রবেশ} Cloud user as a Service (CuaS) নাম দিয়ে ব্রাউজারসহ কম্পিউটার, ট্যাবলেট বা ফোন থেকে ক্লাউডে প্রবেশের ক্লায়েন্ট অংশ বোঝানো হয়। Chrome/Firefox প্রবেশের টুল; নিজেরা ক্লাউড স্টোরেজ নয়। এটি NIST-এর চতুর্থ সেবা মডেল নয়। প্রধান তিনটি মডেল IaaS, PaaS ও SaaS।}

\subsection{Cloud Storage, Sync ও পুনরুদ্ধার}
\label{ch2:c10-11}
নেটওয়ার্ক দিয়ে ক্লাউড সার্ভার বা ডেটা সেন্টারে ফাইল ও ডেটা সংরক্ষণ এবং ব্যবহারের ব্যবস্থাকে ক্লাউড স্টোরেজ বলে। উপযুক্ত প্রবেশাধিকার থাকলে বিভিন্ন ডিভাইস থেকে ফাইল, ছবি ও ভিডিও ব্যবহার, ভাগাভাগি এবং যৌথভাবে কাজ করা যায়। Google Drive, OneDrive, Dropbox ও iCloud পরিচিত উদাহরণ।
\chTwoDiagramRow{cloudbackup}{পরিবর্তন ছড়ানো ও ব্যাকআপের পার্থক্য}{Sync একই পরিবর্তন সংযুক্ত ডিভাইস ও ক্লাউডে ছড়িয়ে দেয়। ভুল করে ফাইল মুছে ফেললে সেই পরিবর্তনও ছড়াতে পারে। ব্যাকআপ পুনরুদ্ধারের জন্য রাখা কপি। আগের সংস্করণের ইতিহাস, পৃথক কপি ও পুনরুদ্ধারের পরীক্ষা প্রয়োজন। সেবাদাতার সার্ভার সচল থাকা এবং নিজের মুছে ফেলা ফাইল ফেরত পাওয়া আলাদা বিষয়।}

\subsection{সুবিধা, খরচ ও সীমাবদ্ধতা}
\label{ch2:c10-12}
\begin{tabularx}{\linewidth}{|p{.28\linewidth}|X|X|}\hline
\textbf{বিষয়} & \textbf{প্রাসঙ্গিক সুবিধা} & \textbf{সীমা/খরচ}\\\hline
সম্পদ ও রক্ষণাবেক্ষণ & নিজের সব ভৌত সার্ভার কিনতে হয় না; অনেক অংশ সেবাদাতা পরিচালনা করে & সেবাদাতার ওপর নির্ভরশীলতা; মডেল অনুযায়ী নিয়ন্ত্রণ ও নিজস্ব দায়িত্ব বদলায়\\\hline
চাহিদা ও খরচ & দ্রুত সম্পদ বাড়ানো বা কমানো; প্রাথমিক বিনিয়োগ কমতে পারে & দীর্ঘদিন বেশি সংরক্ষণ, স্থানান্তর ও প্রক্রিয়াকরণে বিল বাড়ে\\\hline
প্রবেশ ও যৌথ কাজ & উপযুক্ত নেটওয়ার্কে বিভিন্ন ডিভাইস থেকে একই ডকুমেন্টে কাজ & নেটওয়ার্ক বা সেবাদাতার সেবা বন্ধ থাকলে কাজ ব্যাহত; সব কাজ অফলাইনে করা যায় না\\\hline
ডেটা ব্যবস্থাপনা & ভাগাভাগি ও পুনরুদ্ধারের ব্যবস্থা সহজ করা যায় & গোপনীয়তা, ভুল করে মুছে ফেলা ও অননুমোদিত ভাগাভাগির ঝুঁকি; আগের সংস্করণ বা ব্যাকআপ দরকার\\\hline
পরিবর্তন ও স্থানান্তর & নতুন সম্পদ দ্রুত নেওয়া যায় & সেবাদাতা বদলাতে ডেটা ও টুল সরানো জটিল হতে পারে; Vendor Lock-in\\\hline
\end{tabularx}
স্থানীয় ব্যবস্থায় প্রতিষ্ঠান হার্ডওয়্যার, সক্ষমতা ও রক্ষণাবেক্ষণ নিজের মতো পরিচালনা করে। সক্ষমতা বাড়াতে নতুন যন্ত্র ও সময় প্রয়োজন। ক্লাউডে সম্পদ দ্রুত বরাদ্দ করা গেলেও বিল, সংযোগ, সেবাদাতার শর্ত এবং ডেটা রাখার নীতি বিবেচ্য। Pay-as-you-go সবসময় কম মোট খরচ নিশ্চিত করে না। ব্যবহার ও বিল মাপতে হয় এবং অপ্রয়োজনীয় সম্পদ বন্ধ রাখতে হয়।

\subsection{Security ও Shared Responsibility}
\label{ch2:c10-13}
সেবাদাতা ক্লাউড অবকাঠামোর নিরাপত্তা ও পরিচালনার দায়িত্ব নেয়। ব্যবহারকারীর অ্যাকাউন্ট, ডেটা ও প্রবেশাধিকারের দায়িত্ব থাকে। মডেল অনুযায়ী OS ও অ্যাপ্লিকেশন পরিচালনার দায়িত্ব অংশ \ref{ch2:c10-03}-এ আছে। সেবাদাতা আছে বলেই ব্যবহারকারীর সব দায়িত্ব শেষ হয় না।
\chTwoPair{\textbf{কীভাবে সুরক্ষা বাড়তে পারে?}\par
সেবাদাতার প্রবেশাধিকার নিয়ন্ত্রণ, নিয়মিত অবকাঠামো হালনাগাদ ও বিশেষজ্ঞ তদারকি তথ্য রক্ষায় সাহায্য করে। ব্যবহারকারীর শক্তিশালী পাসওয়ার্ড, Multi-factor Authentication (MFA), সীমিত প্রবেশাধিকার, সতর্ক ভাগাভাগি ও উপযুক্ত এনক্রিপশনও প্রয়োজন। MFA-তে ভিন্ন উপাদান দিয়ে অ্যাকাউন্ট যাচাই করা হয়। যেমন, পাসওয়ার্ডের সঙ্গে নিজের ডিভাইসের অনুমোদন।}{\textbf{ঝুঁকি কোথায় থাকে?}\par
অ্যাকাউন্ট চুরি, উন্মুক্ত শেয়ার লিংক, ভুল অনুমতি বা সেবাদাতার ত্রুটিতে অননুমোদিত প্রবেশ ঘটতে পারে। অন্যের সেবায় সংবেদনশীল ডেটা রাখার আগে সংরক্ষণ ও ভাগাভাগির নীতি এবং নিয়ন্ত্রণ দেখতে হয়। অনিরাপদ নেটওয়ার্কে অ্যাকাউন্ট ব্যবহারে সতর্কতা প্রয়োজন। এনক্রিপশন হারানো ফাইল ফেরানোর ব্যবস্থা নয়। ব্যাকআপ ও পুনরুদ্ধারের আলোচনা অংশ \ref{ch2:c10-11}-এ আছে।}

\subsection{Cloud-এর ধারণা, ইতিহাস ও ব্যবহার মাপা}
\chTwoPair{\textbf{ধারণা ও সেবার বিকাশ}\par
John McCarthy ১৯৬১ সালে কম্পিউটিংকে ইউটিলিটি সেবা হিসেবে সংগঠিত করার সম্ভাবনা আলোচনা করেন। পরে নেটওয়ার্ক, ভার্চুয়ালাইজেশন ও ডেটা সেন্টারের উন্নয়নে আধুনিক ক্লাউড গড়ে ওঠে। ২০০৬ সালে Amazon EC2-এর পরীক্ষামূলক Beta সেবা গুরুত্বপূর্ণ বাণিজ্যিক ধাপ। McCarthy-র ধারণা, EC2 চালু ও NIST-এর সংজ্ঞা পৃথক ঘটনা। এগুলো পুরো ক্লাউডের একক যন্ত্র আবিষ্কারের সাল নয়।}{\textbf{একটি কাল্পনিক খরচের হিসাব}\par
স্টোরেজের মাসিক খরচ প্রতি GB ০.৫০ টাকা এবং প্রক্রিয়াকরণের খরচ প্রতি ঘণ্টা ২ টাকা ধরা হলো। ১০০ GB স্টোরেজ ও ১২০ Compute-hour-এর খরচ $100\times0.50+120\times2=290$ টাকা। এটি শেখানোর কাল্পনিক উদাহরণ; সেবাদাতার বর্তমান মূল্য নয়। ডেটা স্থানান্তর, কর ও অন্যান্য খরচ থাকলে পৃথকভাবে যোগ করতে হবে।}
ক্লাউডের জন্য নির্দিষ্ট সর্বোচ্চ দূরত্ব নেই। ক্লায়েন্টের নেটওয়ার্ক সংযোগ ও সেবা পাওয়ার সুযোগ প্রয়োজন। GB সংরক্ষণের পরিমাণ, Mbps ডেটা স্থানান্তরের হার এবং Compute-hour প্রক্রিয়াকরণের ব্যবহৃত সময় বোঝায়। তিনটির একক ও বিল আলাদা। অ্যাকাউন্টসংখ্যা বা বিনামূল্যের স্টোরেজের সীমা সেবাদাতার শর্ত অনুযায়ী বদলায়। একটি সংখ্যাকে ক্লাউডের সংজ্ঞার অংশ ধরা যায় না।

\subsection{কলেজের ভর্তি, ক্লাস ও Result-এর প্রয়োজন}
\label{ch2:c10-14}
\begin{examplebox}{একটি বাস্তব সমস্যা থেকে সেবা নির্বাচন}
কলেজে ভর্তি চলাকালে ওয়েবসাইটে ব্যবহারকারীর চাপ বাড়ে। শিক্ষকরা নোট ও অ্যাসাইনমেন্ট ভাগাভাগি করেন। ফলাফলে প্রবেশাধিকার সীমিত রাখা প্রয়োজন। নোট, ই-মেইল, অনলাইন ডকুমেন্ট, উপস্থিতি ও সভার জন্য প্রস্তুত SaaS ব্যবহার করা যায়। নিজস্ব ভর্তি অ্যাপ্লিকেশন তৈরিতে পরিচালিত প্ল্যাটফর্ম বা PaaS সুবিধাজনক। নির্দিষ্ট OS ও সার্ভারের নিয়ন্ত্রণ প্রয়োজন হলে IaaS বিবেচনা করা যায়।

চাহিদা বাড়লে সম্পদ বাড়ানো, পরে কমানো এবং ব্যবহার ও বিল পর্যবেক্ষণ অংশ \ref{ch2:c10-02}-এর বৈশিষ্ট্যের প্রয়োগ। ফলাফলের জন্য প্রতিষ্ঠানের নিজস্ব Private ব্যবস্থা বিবেচনা করা যায়। Public Admission Cloud-এর সঙ্গে অনুমোদিত ডেটা বিনিময়ের সমন্বিত সংযোগ থাকলে Hybrid হতে পারে। গবেষণাগারের সাময়িক বড় প্রক্রিয়াকরণের কাজেও চাহিদামতো সম্পদ নেওয়া উপকারী। বাস্তব নকশায় উপযুক্ত অনুমতি, ব্যাকআপ ও দায়িত্বের সীমা রাখতে হয়।
\end{examplebox}

\subsection{MCQ অনুশীলন}

\par
\textbf{MCQ}

\begin{multicols}{2}\begin{enumerate}
\item\chTwoMCQBlock{Public Cloud সাধারণত কোন মাধ্যমে ব্যবহার করা হয়?
\par ক. শুধু স্থানীয় প্রিন্টারে \quad খ. ইন্টারনেটের মাধ্যমে \par গ. শুধু অফলাইন ডকুমেন্টে \quad ঘ. কোনো নেটওয়ার্ক ছাড়াই }
\item\chTwoMCQBlock{Google Drive কোন ধরনের সেবার উদাহরণ?
\par ক. ক্লাউড স্টোরেজ \quad খ. স্থানীয় ডিস্কের বিন্যাস \par গ. স্থানীয় ফাইল ডিফ্র্যাগমেন্টেশন \quad ঘ. কিবোর্ডের বিন্যাস }
\item\chTwoMCQBlock{SaaS-এর পূর্ণরূপ কী?
\par ক. Storage as a Service \quad খ. Software as a Service \par গ. System as a Switch \quad ঘ. Server and Scanner }
\item\chTwoMCQBlock{IaaS প্রধানত কী দেয়?
\par ক. প্রস্তুত সফটওয়্যার \quad খ. কম্পিউটিং অবকাঠামো \par গ. কিবোর্ডের বিন্যাস \quad ঘ. মনিটরের উজ্জ্বলতা }
\item\chTwoMCQBlock{অ্যাপ্লিকেশন তৈরি ও চালাতে PaaS প্রধানত কাদের জন্য উপযোগী?
\par ক. সফটওয়্যার উন্নয়নকারী \quad খ. শুধু টাইপিস্ট \par গ. চিত্রশিল্পী \quad ঘ. বিদ্যুৎমিস্ত্রি }
\item\chTwoMCQBlock{Public Cloud সাধারণত কারা পরিচালনা করেন?
\par ক. তৃতীয় পক্ষের সেবাদাতা \quad খ. শুধু একজন শিক্ষার্থী \par গ. একটি স্থানীয় প্রিন্টার \quad ঘ. একটি মাউস }
\item\chTwoMCQBlock{Public ও Private Cloud অবকাঠামো সমন্বিতভাবে যুক্ত হলে কোন মডেল হয়?
\par ক. Hybrid Cloud \quad খ. শুধু Public Cloud \par গ. শুধু Private Cloud \quad ঘ. শুধু Community Cloud }
\item\chTwoMCQBlock{ক্লাউড কম্পিউটিংয়ের একটি সুবিধা কোনটি?
\par ক. প্রয়োজনমতো সম্পদ বাড়ানো বা কমানো \quad খ. শুধু অফলাইনে ব্যবহার \par গ. কেবল পেঁচানো \quad ঘ. মনিটর পরিষ্কার করা }
\item\chTwoMCQBlock{ক্লাউড ব্যবহারে কোন ঝুঁকি বিবেচনা করতে হয়?
\par ক. অননুমোদিত প্রবেশের ঝুঁকি \quad খ. প্রিন্টারের রং \par গ. কিবোর্ডের আকার \quad ঘ. মাউস প্যাড }
\item\chTwoMCQBlock{ক্লাউড অ্যাকাউন্টের নিরাপত্তা বাড়াতে কোন ব্যবস্থা উপযোগী?
\par ক. দুর্বল পাসওয়ার্ড \quad খ. MFA \par গ. পাসওয়ার্ড ভাগাভাগি \quad ঘ. সবার জন্য লগইন উন্মুক্ত রাখা }
\end{enumerate}\end{multicols}

\textbf{উত্তরমালা:}
১. খ, ২. ক, ৩. খ, ৪. খ, ৫. ক, ৬. ক, ৭. ক, ৮. ক, ৯. ক, ১০. খ

\smallskip

\endgroup
